\documentclass[11pt,a4paper]{article}

\usepackage{Act}

\begin{document}
\input{\detokenize{/home/fenarius/Travail/Cours/cpge-info/latex/Macros.tex}}
\ModeExercice
\newcommand{\SPATH}{/home/fenarius/Travail/Cours/cpge-info/docs/mp2i/files/C20/}


\psset{treesep=0.5cm,levelsep=0.8cm}

\section{Inégalités sur les arbres rouge-noirs}

On veut montrer que pour tout arbre rouge-noir $t$ :
\begin{itemize}
	\item $h(t) \leqslant 2b(t)$
	\item $2^{b(t)} \leqslant n+1$
\end{itemize}

On procède par induction structurelle sur $t$ :
\begin{itemize}
	\item Si $t$ est l'arbre vide (seul axiome), alors $h(t)=-1$, $n=0$, $b(t) = 0$ et les inégalités sont vérifiées.
	\item Montrons la conservation par application de l'unique règle d'inférence. On considère un arbre rouge noire $t=(g, r, d)$ et on suppose les deux inégalités vérifiées pour $g$ et $d$. Alors :
	      \begin{itemize}
		      \item Si $r$ est noire, on a alors : d'une part $b(t) = b(l)+ 1$ et $b(t) = b(g) +1$. Par définition de la hauteur : \\
		            $\begin{array}{llll}
				            h(t) & =         & 1+\max(h(g),h(d)) & \text{ par hypothèse d'induction, } h(g) \leqslant 2b(g) \text{ et }  h(d) \leqslant 2b(d) \text{ donc, } \\
				            h(t) & \leqslant & 1 + 2(b(t) -1)    &                                                                                                           \\
				            h(t) & \leqslant & 2b(t) -1          &                                                                                                           \\
				            h(t) & \leqslant & 2b(t)             &                                                                                                           \\
			            \end{array}$
		            Par définition de la taille d'un arbre binaire : \\
		            $\begin{array}{llll}
				            n(t) + 1 & =         & n(g) + n(d) + 2         & \\
				            n(t) + 1 & =         & (n(g)+1) + (n(d)+1)     & \text{ par hypothèse d'induction }\\
				            n(t) + 1 & \geqslant & 2^{b(g)} + 2^{b(d)}     &\\
				            n(t) + 1 & \geqslant & 2^{b(t)-1} + 2^{b(t)-1} & \text{ car la racine est noire}\\
				            n(t) + 1 & \geqslant & 2^{b(t)}                &\\
			            \end{array}$
            \item Si $r$ est rouge, alors $b(t) = b(g) = b(d)$, le cas où l'arbre contient un seul noeud est trivial. Sinon, $t$ a nécessairement deux sous arbres \textit{non vides} sinon la hauteur noire serait différente et ces deux sous arbres ont une racine noire (car leur père est rouge). On note alors $gg$, $gr$, $rg$ et $dd$ les sous arbres des deux sous arbres $g$ et $d$. \\
            $\begin{array}{llll}
                h(t) & = & 1 + \max(h(g),h(d)) & \\
                h(t) & = & 2 + \max(h(gg),h(gd),h(dg),h(dd)) & \text{ or ces 4 arbres ontmême hauteur noire : } b(t)-1\\
                h(t) & \leqslant & 2 + 2 (b(t)-1) & \\
                h(t) & \leqslant & 2b(t) & \\
            \end{array}$
            Par définition de la taille d'un arbre binaire, \\
                $\begin{array}{llll}
				    n(t) + 1 & =         & n(g) + n(d) + 2         & \\
                    n(t) + 1 & =         & (n(g)+1) + (n(d)+1)     & \text{ par hypothèse d'induction }\\
				    n(t) + 1 & \geqslant & 2^{b(g)} + 2^{b(d)}     & \text{ or } b(t) = b(g) = b(d)\\ 
                    n(t) + 1 & \geqslant & 2^{b(t)+1}     &\\
                    n(t) + 1 & \geqslant & 2^{b(t)}     &\\
                \end{array}$
	      \end{itemize}
\end{itemize}

\end{document}